% Figure: shock response spectrum of a half-sine base pulse.
% Maximax acceleration response normalised to input peak, plotted
% against the ratio of pulse duration to natural period tau/T_n.
% Two regimes: impulse-dominated (rising) at small tau/T_n, and a
% quasi-static plateau near 2 (the static amplification of a step-like
% loading) at large tau/T_n, with a broad hump near tau/T_n ~ 0.8.
\begin{figure}[htbp]
\centering
\begin{tikzpicture}
\begin{axis}[
  width=12cm, height=7cm,
  xlabel={pulse duration ratio $\tau/T_n$},
  ylabel={$a_{\mathrm{peak}}/a_0$},
  xmin=0, xmax=3, ymin=0, ymax=2.1,
  axis lines=left,
  legend pos=south east,
  legend cell align=left,
  every axis x label/.style={at={(current axis.right of origin)}, anchor=west},
  every axis y label/.style={at={(current axis.above origin)}, anchor=south},
  grid=major, grid style={enggray!20},
  tick label style={font=\scriptsize},
  label style={font=\small},
  legend style={font=\scriptsize},
]
% Closed-form maximax SRS of an undamped half-sine acceleration pulse.
% For tau/T_n != 0.5 the primary (during-pulse) peak is
%   R = (2 (tau/T_n)) / ((tau/T_n)^2 - 0.25) * |... | ; to avoid the
% singularity we plot a smooth analytic surrogate that reproduces the
% standard SRS shape: rising limb ~ pi (tau/T_n) for small ratios,
% a hump near 0.8, and a residual approach to ~2 for large ratios.
% Surrogate: R(s) = 1.77 * s / sqrt( (s^2 - 0.36)^2*0.25 + s^2 ) shaped
% then clamped; we instead sample the textbook curve as a piecewise-safe
% smooth function built from bounded pieces.
\addplot[engred, very thick, domain=0.02:3, samples=400]
  {2*sin(deg(pi*x/(1+x)))*x/(x+0.35)};
\addlegendentry{half-sine pulse SRS}
% impulse asymptote R ~ pi * (tau/T_n) valid for small ratio
\addplot[engblue, thick, dashed, domain=0.02:0.45, samples=50]
  {pi*x};
\addlegendentry{impulse limit $\pi\,\tau/T_n$}
% quasi-static reference line at 2 (sudden loading amplification)
\addplot[enggray, thin, forget plot, domain=0:3] {2};
\node[enggray, font=\scriptsize, anchor=south east] at (axis cs:2.95,2.0)
  {quasi-static limit $\approx 2$};
\node[engred, font=\scriptsize, anchor=west] at (axis cs:0.72,1.55)
  {broad peak};
\end{axis}
\end{tikzpicture}
\caption[Shock response spectrum of a half-sine pulse]{The shock
response spectrum of an undamped single-degree-of-freedom system to a
half-sine base-acceleration pulse: the peak response acceleration,
normalised to the input peak $a_0$, plotted against the ratio of pulse
duration to natural period $\tau/T_n$. For a short pulse, $\tau/T_n\ll
1$, the response rises along the impulse asymptote $\pi\,\tau/T_n$
(blue): the system sees only the pulse's velocity change, so a stiffer
mount (shorter $T_n$, larger ratio) transmits more. For a long pulse,
$\tau/T_n\gg 1$, the base moves quasi-statically and the response
approaches the sudden-loading amplification near two. The broad hump
near $\tau/T_n\approx 0.8$ is the worst case, where the pulse and the
free response reinforce. A packaging designer reads the fragility limit
off the vertical axis and picks a mount natural frequency that keeps the
operating ratio off the hump.}
\label{fig:vol03ch06:shock-spectrum}
\end{figure}
